\chapter{Conception, Modélisation et Gestion de Projet}

\section{Introduction}

Après la phase de cadrage et de définition des besoins, la conception du projet \textbf{My Smart Budget} a constitué une étape charnière avant le lancement des développements. Cette phase a permis de structurer l'architecture logicielle, de modéliser les données et de définir les processus de collaboration.

L'objectif n'était pas seulement de produire une documentation théorique, mais de bâtir les fondations d'une application \textbf{robuste, évolutive et maintenable}. Pour ce faire, j'ai combiné une approche de modélisation standardisée (UML) avec une méthodologie de gestion de projet Agile (Scrum/Kanban) rigoureuse.

Concrètement, cette phase m'a évité beaucoup d'allers-retours et de refactorisations pendant le développement.

\section{Conception technique et Modélisation UML}

\subsection{Objectifs de la modélisation}

L'utilisation du langage UML (Unified Modeling Language) répond à quatre besoins précis :
\begin{itemize}
    \item \textbf{Structurer} les données et les flux avant l'implémentation.
    \item \textbf{Anticiper} les dépendances complexes entre les modules.
    \item \textbf{Aligner} le code sur les besoins métier (Domain Driven Design).
    \item \textbf{Faciliter} la maintenance future grâce à une documentation visuelle.
\end{itemize}

\subsection{Diagramme de Cas d'Utilisation Global}

Le diagramme ci-dessous synthétise les interactions entre l'utilisateur et le système. Il met en évidence les périmètres fonctionnels clés : gestion de compte, opérations financières et pilotage budgétaire.

\begin{figure}[H]
    \centering
    \safefig{chapters/assets/images/cas-utilisation.png}{0.90\textwidth}{cas-utilisation.png}
    \caption{Diagramme de Cas d'Utilisation global — My Smart Budget}
    \label{fig:uml_global}
\end{figure}

\subsection{Modèle Conceptuel de Données (Détail des Entités)}

La structure des données repose sur trois modules interconnectés :

\begin{description}
    \item[1. Module Utilisateur] \hfill \\
    Gère l'identité et les préférences.
    \begin{itemize}
        \item \textbf{User :} Email unique, mot de passe haché (Bcrypt), date d'inscription.
        \item \textbf{Profile :} Préférences (Devise, Langue), seuils d'alerte personnalisés.
        \item \textbf{Session :} Token JWT (Access + Refresh), expiration, IP.
    \end{itemize}
    
    \item[2. Module Gestion Financière] \hfill \\
    Cœur du système transactionnel.
    \begin{itemize}
        \item \textbf{Transaction :} Montant, date, type (Débit/Crédit), commentaire.
        \item \textbf{Category :} Tag de classification (ex: Alimentation, Loyer) avec icône et couleur.
        \item \textbf{Budget :} Enveloppe limite définie sur une période et une catégorie.
    \end{itemize}

    \item[3. Module Suivi et Alertes] \hfill \\
    Logique proactive.
    \begin{itemize}
        \item \textbf{Alert :} Notification générée automatiquement aux seuils (70\%, 90\%, 100\%).
        \item \textbf{Report :} Agrégat de données pour exports PDF/CSV.
    \end{itemize}
\end{description}

En utilisant UML, j'ai pu voir rapidement si j'avais oublié un cas d'usage important avant d'écrire une ligne de code.

\section{Architecture Technique}

Le projet implémente le pattern architectural \textbf{MVC (Modèle-Vue-Contrôleur)} adapté au web moderne (API RESTful), garantissant une séparation stricte des responsabilités.

\begin{figure}[H]
    \centering
    \safefig{chapters/assets/images/diagramme-packages.png}{0.90\textwidth}{diagramme-packages.png}
    \caption{Architecture technique : Stack Next.js/Express/PostgreSQL+MongoDB et déploiement Docker}
    \label{fig:architecture_tech}
\end{figure}

\begin{tcolorbox}[colback=blue!5!white,colframe=blue!60!black,title=\textbf{Transition Conception \textrightarrow{} Développement}]
La modélisation a servi de guide direct pour l'implémentation :
\begin{itemize}
    \item Les \textbf{entités relationnelles} (User, Transaction, Budget) sont devenues des \textbf{tables PostgreSQL} (via pg-promise).
    \item Les \textbf{données non-structurées} (logs, préférences) sont stockées via des \textbf{schémas Mongoose} (MongoDB).
    \item Les \textbf{Cas d'utilisation} ont défini les routes de l'\textbf{API REST Express}.
    \item Les \textbf{Relations} (1..n) ont dicté les jointures SQL côté PostgreSQL et les références \texttt{ObjectId} côté MongoDB.
\end{itemize}
\end{tcolorbox}

\section{Méthodologie Agile et Conduite de Projet}

\subsection{Adaptation de Scrum pour un projet individuel}

J'ai adopté une approche \textbf{Scrum itérative}, organisée en \textbf{Sprints de 2 semaines}. Cette méthode permet de livrer de la valeur régulièrement tout en ajustant les priorités en fonction des retours (auto-feedback ou tests utilisateurs).

\subsection{Rituels Agiles appliqués}

Bien que travaillant seul, je me suis astreint à des rituels stricts pour maintenir le rythme. En pratique, ces rituels m'ont aidé à garder un rythme régulier et à ne pas laisser certaines tâches traîner :

\begin{itemize}
    \item \textbf{Daily Meeting (10 min)} : Chaque matin, revue des tâches de la veille et définition des objectifs du jour.
    \item \textbf{Sprint Planning (1h / 14j)} : Sélection des User Stories du backlog et découpage en tâches techniques.
    \item \textbf{Sprint Review (30 min / 14j)} : Test fonctionnel des livrables ("Is it working?").
    \item \textbf{Sprint Retrospective (30 min / 14j)} : Analyse du processus ("How to work better?").
\end{itemize}

\begin{tcolorbox}[colback=yellow!5!white,colframe=yellow!60!black,title=\textbf{Exemple concret : Sprint du 14/02/2026}]
\textbf{Lors de la Sprint Review :}
J'ai présenté la fonctionnalité \textit{"CRUD Transactions"}. Un bug a été détecté sur la route DELETE (utilisation de \texttt{db.one} au lieu de \texttt{db.none} sur une requête sans retour).
\begin{itemize}
    \item[$\rightarrow$] \textbf{Action :} Correction immédiate et fermeture de l'issue \#15 \textit{"US-03 – CRUD Transactions catégorisées"} avec commentaire de résolution.
\end{itemize}

\textbf{Lors de la Rétrospective :}
J'ai noté des URLs hardcodées (\texttt{localhost:5001}) dans le frontend, incompatibles avec la production.
\begin{itemize}
    \item[$\rightarrow$] \textbf{Action :} Remplacement systématique par \texttt{process.env.NEXT\_PUBLIC\_API\_URL} et déploiement automatique via CI/CD.
\end{itemize}
\end{tcolorbox}

\subsection{Métriques de suivi et Qualité (KPI)}

Le suivi du projet est piloté par la donnée via le fichier \texttt{metrics.md}, mis à jour à chaque fin de sprint.

\begin{center}
\small
\begin{tabular}{|l|p{5cm}|c|}
\hline
\textbf{Indicateur} & \textbf{Objectif} & \textbf{Fréquence} \\
\hline
Vélocité & $\geq 8$ issues / sprint & Bi-hebdo \\
\hline
Qualité CI/CD & 100\% de succès sur \texttt{main} & Hebdo \\
\hline
Stabilité & $< 5\%$ de bugs post-merge & Sprint Review \\
\hline
Lead Time & $< 3$ jours / fonctionnalité & Hebdo \\
\hline
Code Coverage & $> 70\%$ (Backend Jest) & Sprint Review \\
\hline
\end{tabular}
\end{center}

Ces métriques montrent que le projet avance à un rythme stable et que la qualité est suivie dans le temps.

\begin{tcolorbox}[colback=green!5!white,colframe=green!60!black,title=\textbf{Extrait réel du rapport de Sprint POC (02/04/2026)}]
\begin{verbatim}
## Sprint POC (19/03/2026 - 02/04/2026)
- Issues fermées: 8 (#7, #13, #15, #21, #23, #28, #29, #30)
- Lead Time: 2.1 jours/issue [OK]
- CI Success (main): pipeline vert (build + deploy + SonarCloud)
- Bugs Prod: 2 corrigés (CORS www/non-www, JWT_SECRET manquant)
- Déploiement: smarterbudget.net opérationnel en HTTPS
\end{verbatim}
\end{tcolorbox}

\begin{figure}[H]
    \centering
    \safefig{chapters/assets/images/github-velocity.png}{0.85\textwidth}{github-velocity.png}
    \caption{Graphique de vélocité GitHub Projects (Issues clôturées par Sprint)}
    \label{fig:velocity}
\end{figure}

\subsection{Roadmap et Milestones}

Le projet est jalonné par 4 Milestones GitHub, chacune avec des critères d'entrée (DoR) et de sortie (DoD) stricts. Le jalon POC est livré ; les trois suivants constituent la roadmap active.

\begin{center}
\small
\begin{tabular}{|l|c|l|p{4.5cm}|}
\hline
\textbf{Jalon} & \textbf{Date} & \textbf{Statut} & \textbf{Livrable Principal} \\
\hline
\textbf{POC} & 02/04/2026 & \textcolor{green!60!black}{Livré} & Backend Node.js, connexion PostgreSQL + MongoDB, première route API. \\
\hline
\textbf{M1 — MVP (v1.0)} & 16/04/2026 & \textcolor{orange}{En cours} & Auth JWT, CRUD Transactions, Enveloppes budgétaires, Dashboard connecté API. \\
\hline
\textbf{M2 — Bêta (v1.1)} & 30/04/2026 & \textcolor{gray}{Backlog} & Import CSV, Export JSON, Dashboard analytique avancé, SonarCloud $>$ 80\%. \\
\hline
\textbf{v2.0 — Finale} & 15/05/2026 & \textcolor{gray}{Backlog} & Notifications budgétaires (80\%/100\%), Backoffice admin, Documentation finale. \\
\hline
\end{tabular}
\end{center}

\begin{figure}[H]
    \centering
    \safefig{chapters/assets/images/github-milestones.png}{0.85\textwidth}{github-milestones.png}
    \caption{Suivi des Milestones sur GitHub}
    \label{fig:milestones}
\end{figure}

\subsection{Planification et Burndown Chart}

Les estimations reposent sur une complexité relative (Story Points), où \textbf{1 SP $\approx$ 1 jour homme}.
\begin{itemize}
    \item \textbf{Total Projet :} 21 Story Points (estimé).
    \item \textbf{Vélocité moyenne :} 7 SP / sprint.
\end{itemize}

\begin{figure}[H]
    \centering
    \safefig{chapters/assets/images/burndown-sprint3.png}{0.85\textwidth}{burndown-sprint3.png}
    \caption{Burndown Chart du Sprint 3 : visualisation du travail restant jour par jour}
    \label{fig:burndown}
\end{figure}

\section{Traçabilité totale : De l'idée au code}

J'ai mis en place une chaîne de traçabilité complète pour lier le besoin métier (User Story) à la réalisation technique.

\subsection{Le flux de traçabilité (Workflow)}

\begin{enumerate}
    \item \textbf{User Story :} Définition du besoin (ex: US-03 "Créer transaction").
    \item \textbf{Issue GitHub :} Création du ticket technique (\#12).
    \item \textbf{Branche :} Création d'une branche dédiée (\texttt{feature/us-03...}).
    \item \textbf{Pull Request :} Demande de fusion vers \texttt{main} (\#34).
    \item \textbf{CI/CD :} Exécution automatique des tests (\texttt{ci-tests.yml}).
\end{enumerate}

\begin{center}
\small
\begin{tabular}{|p{2.5cm}|p{2cm}|p{4cm}|p{4cm}|}
\hline
\textbf{User Story} & \textbf{Issue} & \textbf{Branche \& PR} & \textbf{Tests \& CI} \\
\hline
US-01 Auth & \#7 & \texttt{feat/auth} $\rightarrow$ PR\#21 & \texttt{auth.spec.js} \\
\hline
US-03 Transac. & \#12 & \texttt{feat/crud} $\rightarrow$ PR\#34 & \texttt{transaction.spec.js} \\
\hline
US-06 CSV & \#24 & \texttt{feat/csv} $\rightarrow$ PR\#55 & \texttt{import.spec.js} \\
\hline
\end{tabular}
\end{center}

\subsection{Preuves visuelles de l'organisation GitHub}

\begin{figure}[H]
    \centering
    \safefig{chapters/assets/images/github-kanban-board.png}{0.85\textwidth}{github-kanban-board.png}
    \caption{Tableau Kanban : vue globale de l'avancement}
    \label{fig:kanban}
\end{figure}

\begin{figure}[H]
    \centering
    \safefig{chapters/assets/images/github-issue-estimation.png}{0.85\textwidth}{github-issue-estimation.png}
    \caption{Issue détaillée avec estimation (Story Points) et critères d'acceptation}
    \label{fig:issue}
\end{figure}

\begin{figure}[H]
    \centering
    \safefig{chapters/assets/images/github-pr.png}{0.85\textwidth}{github-pr.png}
    \caption{Pull Request validée avec checklist et revue de code}
    \label{fig:pr}
\end{figure}

\begin{figure}[H]
    \centering
    \safefig{chapters/assets/images/github-actions-success.png}{0.85\textwidth}{github-actions-success.png}
    \caption{Pipeline CI/CD : tests passés avec succès avant merge}
    \label{fig:ci_cd}
\end{figure}

\section{Conclusion du chapitre}

Cette phase de conception et d'organisation a permis de :
\begin{itemize}
    \item \textbf{Sécuriser le développement} grâce à une architecture validée (UML/MVC).
    \item \textbf{Garantir la qualité} par l'application rigoureuse des processus Agiles.
    \item \textbf{Prouver le professionnalisme} via une traçabilité complète sur GitHub.
\end{itemize}

Les fondations étant posées, le chapitre suivant détaillera la réalisation technique et le développement de l'application.